\section{Method}
\subsection{Problem Formulation}
Given a telemetry window $x_{1:T}$ and optional context $c$, the model predicts a health state $y \in \{\text{normal}, \text{warning}, \text{critical}\}$, an optional fault or subsystem label, and a maintenance action. The architecture is designed so that the same latent reasoning block can be applied repeatedly before producing a final decision.

\subsection{Your Solution}
The proposed solution is an adaptive looped reasoning model for predictive maintenance. Instead of making a single forward-pass decision, the model encodes a telemetry window into a latent state, repeatedly refines that state with a shared recurrent block, and uses a learned exit rule to decide when enough internal computation has been used. This is paired with multitask heads for health state, maintenance action, and subsystem attribution so the system produces operationally meaningful outputs rather than only a class label.

\subsection{Architecture}
The model has three stages. A window encoder maps raw sensor inputs to a latent representation. A shared loop block refines that latent state through repeated updates. A confidence-based exit policy decides whether additional internal steps are needed. The final state feeds task-specific heads for classification and action selection. The design borrows the core idea of adaptive halting from adaptive computation time and the repeated refinement pattern from latent recurrence models \cite{Graves2016ACT,Dehghani2018UniversalTransformers,Banino2021PonderNet}.

\begin{equation}
h_0 = f_{\text{enc}}(x_{1:T}, c)
\end{equation}
\begin{equation}
h_{k+1} = f_{\theta}(h_k, c)
\end{equation}
\begin{equation}
p_{\text{exit}}^{(k)} = \max \left(\text{softmax}(g(h_k))\right)
\end{equation}

The loop is unrolled until the confidence threshold is satisfied or a maximum depth is reached. This structure makes compute usage data-dependent rather than fixed, which is central to the main hypothesis. The looped reasoning framing is also aligned with the current generation of latent reasoning models, which treat repeated internal computation as a mechanism for stronger knowledge manipulation rather than simply larger parameter counts \cite{Zhu2025Ouro}.

\subsection{Baselines}
The evaluation compares the adaptive looped model against three baselines: a direct GRU-based classifier, a fixed-depth version of the same looped architecture, and a pretrained text-transformer baseline. The LLM baseline serializes each telemetry window into a compact textual summary and applies a frozen DistilBERT encoder with lightweight multitask heads. Together, these comparisons isolate the value of recurrence, adaptive early exit, and numeric telemetry modeling versus text-serialized transformer inference.

\subsection{Implementation Notes}
The implementation uses a GRU encoder with hidden size 128, AdamW optimization, and a maximum loop depth of six steps. The adaptive model halts when the classifier confidence exceeds 0.8. The pipeline records average depth per example, latency benchmarks, validation and test confusion matrices, and full metric reports for each run.
